
\documentclass{article}
\usepackage{colortbl}
\usepackage{makecell}
\usepackage{multirow}
\usepackage{supertabular}

\begin{document}

\newcounter{utterance}

\centering \large Interaction Transcript for game `cladder', experiment `full\_v1.5\_default', episode 2679 with Qwen3.6{-}27B.
\vspace{24pt}

{ \footnotesize  \setcounter{utterance}{1}
\setlength{\tabcolsep}{0pt}
\begin{supertabular}{c@{$\;$}|p{.15\linewidth}@{}p{.15\linewidth}p{.15\linewidth}p{.15\linewidth}p{.15\linewidth}p{.15\linewidth}}
    \# & \multicolumn{2}{c}{Player} && \multicolumn{2}{c}{Game Master} \\
    \hline

    \theutterance \stepcounter{utterance}  
    & & & \multicolumn{4}{p{0.6\linewidth}}{
        \cellcolor[rgb]{0.9,0.9,0.9}{
            \makecell[{{p{\linewidth}}}]{
                \texttt{\tiny{[P1$\langle$GM]}}
                \texttt{You are an expert in causal inference. The following question is not a typical commonsense query, but rather a meticulously designed question created by a professor specializing in causal inference, intended to assess the students' mastery of the course content.} \\
\texttt{Imagine a self{-}contained, hypothetical world with only the following conditions, and without any unmentioned factors or causal relationships: Gene has a direct effect on smoking and lung cancer. Pollution has a direct effect on smoking and lung cancer. Smoking has a direct effect on lung cancer. Pollution is unobserved. For people with nonsmoking genes and nonsmokers, the probability of lung cancer is 56\%. For people with nonsmoking genes and smokers, the probability of lung cancer is 42\%. For people with smoking genes and nonsmokers, the probability of lung cancer is 60\%. For people with smoking genes and smokers, the probability of lung cancer is 40\%. For people with nonsmoking genes and with low pollution, the probability of smoking is 76\%. For people with nonsmoking genes and with high pollution, the probability of smoking is 48\%. For people with smoking genes and with low pollution, the probability of smoking is 46\%. For people with smoking genes and with high pollution, the probability of smoking is 21\%. The overall probability of high pollution is 41\%. Does gene positively affect lung cancer through smoking?} \\
\texttt{Start your answer with 'Yes' or 'No', followed by additional reasoning or evidence to support your explanation.} \\
            }
        }
    }
     \\ \\

    \theutterance \stepcounter{utterance}  
    & \multicolumn{4}{p{0.6\linewidth}}{
        \cellcolor[rgb]{0.561,0.737,0.561}{
            \makecell[{{p{\linewidth}}}]{
                \texttt{\tiny{[P1$\rangle$GM]}}
                \texttt{No} \\
\\ 
\texttt{To determine whether Gene positively affects Lung Cancer through Smoking, we must calculate the **Causal Mediation Effect** (specifically, the Natural Direct Effect or the Controlled Direct Effect, but in this context, we are looking for the sign of the indirect effect or the total effect mediated by smoking). However, the question asks if Gene *positively* affects Lung Cancer *through* Smoking. This implies we need to check the sign of the indirect effect: \$Gene \textbackslash{}rightarrow Smoking \textbackslash{}rightarrow Lung Cancer\$.} \\
\\ 
\texttt{Let's break down the causal structure and the data provided.} \\
\\ 
\texttt{**1. Causal Structure:**} \\
\texttt{*   \$G\$ (Gene) \$\textbackslash{}rightarrow\$ \$S\$ (Smoking)} \\
\texttt{*   \$G\$ (Gene) \$\textbackslash{}rightarrow\$ \$L\$ (Lung Cancer)} \\
\texttt{*   \$P\$ (Pollution) \$\textbackslash{}rightarrow\$ \$S\$ (Smoking)} \\
\texttt{*   \$P\$ (Pollution) \$\textbackslash{}rightarrow\$ \$L\$ (Lung Cancer)} \\
\texttt{*   \$S\$ (Smoking) \$\textbackslash{}rightarrow\$ \$L\$ (Lung Cancer)} \\
\texttt{*   \$P\$ is unobserved (confounder for \$S \textbackslash{}rightarrow L\$ and \$G \textbackslash{}rightarrow S\$? No, \$G\$ and \$P\$ are independent unless stated otherwise. The problem states "without any unmentioned factors or causal relationships", implying \$G\$ and \$P\$ are independent root nodes).} \\
\\ 
\texttt{**2. Data Provided} \\
            }
        }
    }
    & & \\ \\

    \theutterance \stepcounter{utterance}  
    & & & \multicolumn{2}{p{0.3\linewidth}}{
        \cellcolor[rgb]{0.9,0.9,0.9}{
            \makecell[{{p{\linewidth}}}]{
                \texttt{\tiny{[GM$|$GM]}}
                \texttt{No} \\
            }
        }
    }
    & & \\ \\

    \theutterance \stepcounter{utterance}  
    & & & \multicolumn{2}{p{0.3\linewidth}}{
        \cellcolor[rgb]{0.9,0.9,0.9}{
            \makecell[{{p{\linewidth}}}]{
                \texttt{\tiny{[GM$|$GM]}}
                \texttt{yes} \\
            }
        }
    }
    & & \\ \\

    \theutterance \stepcounter{utterance}  
    & & & \multicolumn{2}{p{0.3\linewidth}}{
        \cellcolor[rgb]{0.9,0.9,0.9}{
            \makecell[{{p{\linewidth}}}]{
                \texttt{\tiny{[GM$|$GM]}}
                \texttt{game\_result = LOSE} \\
            }
        }
    }
    & & \\ \\

\end{supertabular}
}

\end{document}
